\chapter{Experiments and results}
\label{ch:exp}

\section{Hardware}
\label{sec:hardware}

Every trial in this chapter, from the hardware validation to the simulated user study, runs on the same robot, driven through one of two devices from the same haptic-interface family. This section introduces the physical hardware once, at a high level; the quantitative envelope of the study's own device is tabulated in Appendix~\ref{app:config}.

\subsection{PAL TRIAGo}
\label{sec:hwtriago}

The robot is a PAL Robotics TRIAGo, a bimanual mobile manipulator built around a wheeled base, a lift torso, two 7-DOF arms ending in two-finger grippers, and a 7-DOF sensorised head carrying an RGB-D camera (Figure~\ref{fig:hwtriago}). Every mechanism this thesis controls --- the safety filter of Chapter~\ref{ch:method}, the shared-autonomy arbitration, and the active-vision head servo --- runs on this same kinematic chain, in simulation throughout the study and on the physical unit shown here for the two hardware trials reported in this chapter.

\begin{figure}[t]
\centering
\includegraphics[height=6.8cm]{triago_lateral.png}\hspace{1.5cm}%
\includegraphics[height=6.8cm]{triago_front.png}
\caption{The physical PAL TRIAGo: a lateral view (left), and a front view in the home configuration the robot assumes at power-on, with the two-finger grippers mounted (right).}
\label{fig:hwtriago}
\end{figure}

\subsection{Haption Desktop 6D Compact}
\label{sec:hwdesktop}

The operator's interface for every session of the simulated study is a Haption Desktop 6D Compact (Figure~\ref{fig:haptiondevices}, left), a benchtop, six-degree-of-freedom impedance-type haptic device: it measures the handle's own pose and renders a commanded wrench back onto it, the mechanism both control mappings of Section~\ref{sec:device} are built on. A push-button pad fitted on the handle (Figure~\ref{fig:haptiondevices}, centre) carries the operator's discrete commands; its buttons are programmable, so each can be mapped in software to an arbitrary command (Section~\ref{sec:opcontrols}). Its joint travel and Cartesian workspace are tabulated in Appendix~\ref{app:config}.

\subsection{Haption Virtuose 6D}
\label{sec:hwvirtuose}

The two hardware trials instead use a Haption Virtuose 6D (Figure~\ref{fig:haptiondevices}, right), the full-size member of the same device family: the same six-degree-of-freedom impedance interface and control software, mounted on a floor-standing pedestal with a substantially larger reach than the desktop unit. The two devices share a control interface but not a physical envelope, so the desktop unit's workspace of Appendix~\ref{app:config} does not describe the operator's workspace in the teleoperated hardware trial. Unlike the desktop unit, it also carries a deadman switch on its handle: the device renders forces only while the switch is held, and the hardware trials count the operator as in control only over those intervals.

\begin{figure}[!t]
\centering
\raisebox{-0.5\height}{\includegraphics[height=5.1cm]{haption_desktop.png}}\hfill
\raisebox{-0.5\height}{\includegraphics[width=2.9cm]{button.png}}\hfill
\raisebox{-0.5\height}{\includegraphics[height=5.39cm]{haption_virtuose.png}}
\caption{The haptic hardware used in this thesis: the Desktop 6D Compact (left) for the simulated study, the programmable push-button pad fitted on its handle, through which the operator issues every discrete command (centre), and the full-size Virtuose 6D (right) for the two hardware trials.}
\label{fig:haptiondevices}
\end{figure}

\FloatBarrier

\section{Hardware validation}
\label{sec:hwvalidation}

The perception pipeline is validated separately on the physical robot. The head scans a table carrying two cups; the table plane and both objects are recovered by the RANSAC-based estimator, and the frozen snapshot is loaded as the safety controller's collision model. Teleoperation and barrier-based obstacle avoidance are then exercised against that perceived model, confirming functional end-to-end operation, from depth stream to enforced constraint, but not the geometric accuracy of the perceived model, which is not measured against ground truth on hardware.

\subsection{Start-to-crossed reaching}
\label{sec:hwswap}

The first hardware test targets the controller alone, with vision excluded from the loop. Both arms are driven open-loop, by the same quintic-timed Cartesian reference source used throughout controller development, from a common starting pose --- both grippers raised and held away from the body (Figure~\ref{fig:hwhome}) --- to the pose sampled from the other arm at the start of the motion, so each arm is commanded straight through the space the other one occupies (Figure~\ref{fig:hwcrossed}). That target pair is not jointly reachable without the self-collision and inter-arm rows of the safety filter (Section~\ref{sec:cbf}) becoming active, so the manoeuvre is posed as a conflict between tracking and safety rather than a compliant reach: it tests whether the CLF's tracking slack yields gracefully under that conflict (Section~\ref{sec:clf}), by exactly as much as each arm's adaptively scheduled weight (Section~\ref{sec:sched}) allows, so the controller settles on the best pose the hard constraints admit for both grippers rather than stalling or violating one.

\begin{figure}[t]
\centering
\begin{subfigure}[t]{0.46\linewidth}\centering\includegraphics[width=\linewidth]{TRIAGO_Home_Colored.png}\caption{Starting pose}\label{fig:hwhome}\end{subfigure}\hfill
\begin{subfigure}[t]{0.46\linewidth}\centering\includegraphics[width=\linewidth]{TRIAGO_Crossed.png}\caption{Crossed configuration}\label{fig:hwcrossed}\end{subfigure}
\caption{The physical robot at the two endpoints of the start-to-crossed reaching test: both grippers start raised and away from the body (left) and finish crossed in front of the torso, each having reached the other's starting pose (right).}
\label{fig:hwswap}
\end{figure}

The manoeuvre is logged with the controller's own telemetry (Section~\ref{sec:sched}), reconstructed offline from the recorded rosbag, over a $25.1$~s trial: the shared quintic-timed reference for both arms runs open-loop from $t=0$ to $t\approx10$~s, marked by the dashed vertical line in every panel of Figures~\ref{fig:hwjoints} and~\ref{fig:hwqpdata}, after which the reference is held fixed at the crossed endpoint for the remainder of the recording. The physical robot sustains a mean loop rate of $126$~Hz with $15.7\%$ tick-to-tick jitter against the $150$~Hz nominal --- lower and noisier than in simulation, as expected on hardware --- with the barrier row computed asynchronously and consumed up to two ticks stale, as Section~\ref{sec:concllimitations} describes. This loop-rate telemetry is a window mean over four control ticks, since the telemetry decimation is doubled on hardware: the window-mean tick period has a median of $7.92$~ms, a 95th percentile of $10.6$~ms, a 99th of $11.7$~ms, and a worst case of $13.6$~ms, with $9.61\%$ of windows exceeding the $6.67$~ms nominal period by more than $50\%$ and $0.132\%$ by more than $100\%$; an independent per-tick estimate, from the barrier-row telemetry published every tick, gives a lower median of $7.58$~ms but a heavier tail, up to $15.8$~ms at the 99th percentile and $19.9$~ms worst case. No exact-zero, infeasible-branch command appears in this trial's $761$ decimated command samples, one in four ticks, which is consistent with the branch's rarity but is not a full-resolution count, and a watchdog freeze is not distinguishable from that branch in the recorded data. Low-level execution is faithful throughout: the root-mean-square mismatch between the commanded and the measured joint velocity, across all 14 arm joints, is $5.0\times10^{-3}$~rad/s, so the discrepancies reported below originate in the constrained optimum the QP is solving, not in how well the robot executes it.

\begin{figure}[t]
\hwjointlegend
\begin{subfigure}[t]{0.485\linewidth}\centering\includegraphics{results/hw_q_R.pdf}\caption{Joint positions, right arm}\label{fig:hw-qR}\end{subfigure}\hfill
\begin{subfigure}[t]{0.485\linewidth}\centering\includegraphics{results/hw_q_L.pdf}\caption{Joint positions, left arm}\label{fig:hw-qL}\end{subfigure}
\vspace{0.3em}
\begin{subfigure}[t]{0.485\linewidth}\centering\includegraphics{results/hw_qdot_R.pdf}\caption{Measured joint velocities, right arm}\label{fig:hw-qdR}\end{subfigure}\hfill
\begin{subfigure}[t]{0.485\linewidth}\centering\includegraphics{results/hw_qdot_L.pdf}\caption{Measured joint velocities, left arm}\label{fig:hw-qdL}\end{subfigure}
\caption{Joint-space view of the trial: position and measured velocity of the seven joints of each arm, J1 (shoulder) to J7 (wrist).}
\label{fig:hwjoints}
\end{figure}

Figure~\ref{fig:hwjoints} is the trial seen from the joints. No joint exceeds $0.27$~rad/s at any point, and every velocity trace is smooth: there is no trace of the high-frequency ripple that an unanchored cost otherwise produces on this robot; the measured-velocity rate damping of Section~\ref{sec:rate} suppresses it. The two arms do not stop when the reference does: the joints keep moving after $t\approx10$~s at up to $0.2$~rad/s, slowly and continuously, which is the settling that the shadow prices of Figure~\ref{fig:hwqpdata} account for tick by tick --- the arms are still being moved by the barrier and joint-limit rows, not drifting. The two late joint-limit bursts are visible here as the last inflections of the position traces, around $19$~s on the right arm and $22$~s on the left, after which the configurations settle and the velocities decay to zero.

\begin{figure}[p]
\hwrllegend
\begin{subfigure}[t]{0.485\linewidth}\centering\includegraphics{results/hw_lambda_cbf.pdf}\caption{CBF shadow price}\label{fig:hw-lcbf}\end{subfigure}\hfill
\begin{subfigure}[t]{0.485\linewidth}\centering\includegraphics{results/hw_lambda_jl.pdf}\caption{Joint-limit shadow price}\label{fig:hw-ljl}\end{subfigure}
\vspace{0.4em}
\begin{subfigure}[t]{0.485\linewidth}\centering\includegraphics{results/hw_margin.pdf}\caption{CBF safety margin}\label{fig:hw-margin}\end{subfigure}\hfill
\begin{subfigure}[t]{0.485\linewidth}\centering\includegraphics{results/hw_dmin.pdf}\caption{Minimum inter-geometry distance}\label{fig:hw-dmin}\end{subfigure}
\vspace{0.4em}
\begin{subfigure}[t]{0.485\linewidth}\centering\includegraphics{results/hw_dsafe.pdf}\caption{Velocity-dependent safety margin}\label{fig:hw-dsafe}\end{subfigure}\hfill
\begin{subfigure}[t]{0.485\linewidth}\centering\includegraphics{results/hw_wslack.pdf}\caption{Adaptive slack weight}\label{fig:hw-wslack}\end{subfigure}
\vspace{0.4em}
\begin{subfigure}[t]{0.485\linewidth}\centering\includegraphics{results/hw_gov_ep.pdf}\caption{Limiter: position error}\label{fig:hw-govep}\end{subfigure}\hfill
\begin{subfigure}[t]{0.485\linewidth}\centering\includegraphics{results/hw_gov_theta.pdf}\caption{Limiter: orientation error}\label{fig:hw-govtheta}\end{subfigure}
\caption{QP-CLF-CBF and reference-limiter data of the same trial. The vertical axis of (a) is broken between $180$ and $500$ to keep the single $557$ spike in view. Traces (c) and (d) start at $t \approx 6$~s, when the first collision pair enters the activation radius.}
\label{fig:hwqpdata}
\end{figure}

Figures~\ref{fig:hw-lcbf} to~\ref{fig:hw-dsafe} are the safety filter's own account of the conflict. Both collision shadow prices climb together during the transit; the left arm's spikes to the trial's peak of $557$ near $t=6$~s, and both arms show repeated smaller bursts of $100$--$150$ throughout the hold phase, each a fresh barrier engagement rather than one sustained event. That peak and its neighbour are the only two samples above $180$ in the whole trial, so the vertical axis of Figure~\ref{fig:hw-lcbf} is broken between $180$ and $500$: the spike stays visible in the top strip without compressing the bursts that make up the rest of the trace. The joint-limit shadow price is active mainly during the transit, then returns in two brief late bursts, around $19$~s for the right arm and $22$~s for the left. Figures~\ref{fig:hw-margin} and~\ref{fig:hw-dmin} report what those prices are protecting, and both traces likewise only start around this same $t\approx6$~s rather than at $t=0$: both arms begin raised and away from the body and from each other, so no pair is yet within the activation radius (Section~\ref{sec:cbf}), and neither $h_{\text{soft},X}$ nor the raw minimum distance carries a meaningful value against an empty, inactive pair set. From that point on, the CBF safety margin dips to just below zero at its worst ($-2.6$~mm) and otherwise sits within a few millimetres of it for the whole hold phase, while the raw minimum inter-geometry distance floors at $4.2$~cm over the same stretch --- clearance is barrier-limited, the same pattern the human-subject study exhibits at full scale. The velocity-dependent margin $d_{\text{safe},X}$ of~\eqref{eq:dsafe} (Figure~\ref{fig:hw-dsafe}) tracks each arm's own speed directly: it rises to as much as $5.3$~cm during the fast transit and relaxes back toward its $1.5$~cm static floor $d_0$, dotted, as each arm slows.

Figures~\ref{fig:hw-wslack} to~\ref{fig:hw-govtheta} show the two mechanisms that shape what the QP is even asked to do. The slack weight $w_{\delta,X}$ sits at its ceiling before either arm starts moving, collapses to its floor within the first few seconds of motion, and then cycles back toward the ceiling every time that arm's own shadow price relaxes --- the schedule of~\eqref{eq:slacksched}, now read directly from its own output rather than inferred from the slack it produces. Of the reference limiter's four mechanisms (Algorithm~\ref{alg:governor}), only one is exercised by this manoeuvre: the commanded linear speed and angular speed never approach their configured ceilings at any point (not shown), and neither does the orientation-error norm, whose raw and governed traces coincide throughout in Figure~\ref{fig:hw-govtheta} (dashed under solid). Position-error bounding (Figure~\ref{fig:hw-govep}) is the exception, and it is not a brief event: once the right arm's real pose falls behind its raw reference, the governed error is pinned at a common ceiling of about $0.30$~m on $78\%$ of the trial's ticks, and the left arm's error is clamped intermittently over the same stretch it is still converging. The raw traces of the same panel are the residual each arm is left with: at the end of the trial the right gripper sits $0.42$~m and $14.7^{\circ}$ from its commanded pose, the left $0.11$~m and $28.8^{\circ}$ --- each arm converges on a different coordinate of the same target pose, the residual falling wherever its own slack leaves it. This is the same fact the shadow-price and slack panels already show from the QP's side of the loop, seen here from upstream of it: whenever the safety filter is genuinely fighting the reference, the limiter keeps re-issuing one the CLF can still make progress toward, rather than the raw, receding one.

\subsection{Teleoperated bimanual grasping}
\label{sec:hwteleop}

The second hardware trial closes the loop that Section~\ref{sec:hwswap} leaves open: the operator drives the robot directly, through the shared-autonomy architecture of Chapter~\ref{ch:method} rather than an open-loop reference, against the same two-cup tabletop scene introduced at the start of this section. The condition is \textsc{clutch} full guidance (CFB, Table~\ref{tab:cells}): the operator's tether is rendered as both a guidance-biased force on the handle and, through the blending layer, the arm reference itself. The device driving the real robot is a full-size Haption Virtuose 6D, not the Desktop 6D Compact whose joint travel and workspace are tabulated in Appendix~\ref{app:config} for the simulated study --- the two share a control interface but not a physical envelope, so the operator's available motion range in this trial is not the one those tables describe. The trial is recorded over $253.7$~s at a mean loop rate of $109.6$~Hz with $18.6\%$ tick-to-tick jitter, lower and noisier than the open-loop trial of Section~\ref{sec:hwswap} since the shared-autonomy stack now also carries belief estimation, the grasp state machine, and the blending arbitration on every tick. Its window-mean tick period has a median of $9.22$~ms, a 95th percentile of $12.6$~ms, a 99th of $14.6$~ms, and a worst case of $29.4$~ms, with $34.8\%$ of windows exceeding the $6.67$~ms nominal period by more than $50\%$ and $2.51\%$ by more than $100\%$; the per-tick barrier-row estimate gives a median of $8.97$~ms, a 99th percentile of $18.8$~ms, and a worst case of $36.9$~ms. Neither bag records per-tick QP solve times or in-trial end-to-end command latency; a separate transport measurement on the robot network finds a joint-state delay of mean $3.00$~ms, standard deviation $0.12$~ms, a transport figure rather than a sensor-to-actuation latency. The root-mean-square mismatch between commanded and measured joint velocity across all 14 arm joints is $1.2\times10^{-2}$~rad/s, higher than Section~\ref{sec:hwswap}'s $5.0\times10^{-3}$ but still an order of magnitude below the motions it is tracking, so here too execution fidelity is not the limiting factor in what follows.

% NOTE (not for the reader): radius, table footprint, and the blue cylinder's height in fig:telworld are NOT recovered from this trial's own head-perception snapshot (that pipeline doesn't carry them on any topic this bag records) -- they are retrieved from previously recorded values for the same physical rig, reused here for continuity. The PDF text below deliberately does not say this; it states them as retrieved numbers, not as a fresh physical measurement. Keep this distinction in mind if the provenance of these figures is questioned later.
\textbf{World.} The collision model is the head's own live estimate (Section~\ref{sec:convergence}), not a declared one: the trial launches the camera-driven controller variant, which blocks at startup on the head's converged, latched snapshot and never falls back to a static world file. The world holds two cylinders: the red cylinder's top surface sits at $0.872$~m above the base frame, and Figure~\ref{fig:telworld} gives its radius, the table's own footprint, and the blue cylinder's height. The red cylinder sits on the robot's right and the blue on its left, so throughout this section red is the right arm's object and blue the left's, the same colour convention as the arms themselves.

\begin{figure}[t]
\centering
\includegraphics[width=0.92\linewidth]{results/tel_world_recon.pdf}
\caption{Reconstructed two-cylinder world of the trial: top view (left) and front elevation (right). Candidate points are illustrative, not the recorded cloud.}
\label{fig:telworld}
\end{figure}

\begin{figure}[t]
\centering
\begin{subfigure}[t]{0.30\linewidth}\centering\includegraphics[width=\linewidth]{results/Real_Teleop_10_CFB_SUCCESS/Teleop1_crop.jpeg}\caption{Home configuration}\label{fig:tel-p1}\end{subfigure}\hspace{1.5em}%
\begin{subfigure}[t]{0.30\linewidth}\centering\includegraphics[width=\linewidth]{results/Real_Teleop_10_CFB_SUCCESS/Teleop2_crop.jpeg}\caption{Approaching the near cup}\label{fig:tel-p2}\end{subfigure}
\par\vspace{0.3em}
\begin{subfigure}[t]{0.30\linewidth}\centering\includegraphics[width=\linewidth]{results/Real_Teleop_10_CFB_SUCCESS/Teleop3_crop.jpeg}\caption{Right cup secured}\label{fig:tel-p3}\end{subfigure}\hspace{1.5em}%
\begin{subfigure}[t]{0.30\linewidth}\centering\includegraphics[width=\linewidth]{results/Real_Teleop_10_CFB_SUCCESS/Teleop4_crop.jpeg}\caption{Reaching for the second cup}\label{fig:tel-p4}\end{subfigure}
\caption{The physical trial in sequence: from the home pose, the right arm reaches and secures the first cup, then the left arm engages the second.}
\label{fig:telphotos}
\end{figure}

\textbf{Course of the trial.} Both objects were within reach of either arm; the right arm started active (Figure~\ref{fig:telphotos}). At $t=46.3$~s the operator triggered a grasp on the right cup approached from the side; the state machine held it in alignment for $12.0$~s --- matching the alignment phase's own timeout almost exactly --- without ever issuing a close command, and released it back to the open gripper at $t=58.4$~s: a convergence failure during approach, not a closing or contact failure, since the gripper was never commanded to close at all. The operator repositioned and retried from above at $t=103.2$~s; the gripper closed over $t=109.1$--$111.4$~s and the cup was reported attached at $t=114.4$~s. Control switched to the left arm at $t=132.9$~s, and at $t=171.2$~s the operator triggered a side grasp on the second cup. This attempt was held for $37.8$~s --- longer than either the alignment or the approach phase's own timeout individually --- without a single close command, ending with the gripper explicitly reopened at $t=206.4$~s; the pattern (a long stall at a constant, un-ramped margin, never reaching contact) is the state machine's own signature of a retreat that cannot find clearance, consistent with the wedged or joint-limited case its recovery logic itself checks for. At the trigger, the state machine removed the pair formed by the left gripper and the blue cup from the barrier entirely, rather than relaxing its margin, and held it removed through the stall and the abort retreat that followed; the pair was restored by a fixed $3.0$~s timer at $t=209.4$~s rather than by a measured-clearance test (Section~\ref{sec:grasp}). This revision predates the clearance-gated recovery described there. The controller's own telemetry marks what came next: the CBF shadow price on the left arm, quiet all through the stall, breaks into a ragged, repeatedly re-triggered burst from $t\approx211$ to $217$~s (Figure~\ref{fig:tel-lcbf}) as the arm oscillates against the barrier rather than retreating cleanly, and the operator released the device's deadman switch at $t=217.3$~s, stopping the robot. The recording continues to $253.7$~s, but nothing after the stop reflects the controlled system any longer; every panel below is therefore shown only to $t=230$~s, a few seconds past the stop itself.

\begin{figure}[t]
\hwjointlegend
\begin{subfigure}[t]{0.485\linewidth}\centering\includegraphics{results/tel_q_R.pdf}\caption{Joint positions, right arm}\label{fig:tel-qR}\end{subfigure}\hfill
\begin{subfigure}[t]{0.485\linewidth}\centering\includegraphics{results/tel_q_L.pdf}\caption{Joint positions, left arm}\label{fig:tel-qL}\end{subfigure}
\vspace{0.3em}
\begin{subfigure}[t]{0.485\linewidth}\centering\includegraphics{results/tel_qdot_R.pdf}\caption{Measured joint velocities, right arm}\label{fig:tel-qdR}\end{subfigure}\hfill
\begin{subfigure}[t]{0.485\linewidth}\centering\includegraphics{results/tel_qdot_L.pdf}\caption{Measured joint velocities, left arm}\label{fig:tel-qdL}\end{subfigure}
\caption{Joint-space view of the teleoperated trial, to $t=230$~s (Section~\ref{sec:hwteleop}).}
\label{fig:telqjoints}
\end{figure}

Figure~\ref{fig:telqjoints} shows the joint-space signature of that same course of events: the right arm moves through the approach and grasp, makes one further, sharp joint-6 adjustment at $t\approx127$~s as the switch pins its weight to the rigid, inactive-arm ceiling, and is essentially static thereafter, while the left arm's joint velocities stay near zero until the switch at $t\approx133$~s and then show two distinct episodes, the grasp approach and, sharply, the $t\approx211$--$217$~s window. There the left arm's velocity trace stops being a single smooth motion and becomes a rapid back-and-forth: joint~4 alone reverses direction roughly every $0.08$~s (about $6$~Hz) at up to $0.11$~rad/s, an order of magnitude faster than any direction change asked of either arm earlier in the trial --- the joint-space trace of the oscillation the shadow price already reports.

\begin{figure}[p]
\hwrllegend
\begin{subfigure}[t]{0.485\linewidth}\centering\includegraphics{results/tel_lambda_cbf.pdf}\caption{CBF shadow price}\label{fig:tel-lcbf}\end{subfigure}\hfill
\begin{subfigure}[t]{0.485\linewidth}\centering\includegraphics{results/tel_lambda_jl.pdf}\caption{Joint-limit shadow price}\label{fig:tel-ljl}\end{subfigure}
\vspace{0.4em}
\begin{subfigure}[t]{0.485\linewidth}\centering\includegraphics{results/tel_margin.pdf}\caption{CBF safety margin}\label{fig:tel-margin}\end{subfigure}\hfill
\begin{subfigure}[t]{0.485\linewidth}\centering\includegraphics{results/tel_dmin.pdf}\caption{Minimum inter-geometry distance}\label{fig:tel-dmin}\end{subfigure}
\vspace{0.4em}
\begin{subfigure}[t]{0.485\linewidth}\centering\includegraphics{results/tel_dsafe.pdf}\caption{Velocity-dependent safety margin}\label{fig:tel-dsafe}\end{subfigure}\hfill
\begin{subfigure}[t]{0.485\linewidth}\centering\includegraphics{results/tel_wslack.pdf}\caption{Adaptive slack weight}\label{fig:tel-wslack}\end{subfigure}
\vspace{0.4em}
\begin{subfigure}[t]{0.485\linewidth}\centering\includegraphics{results/tel_gov_ep.pdf}\caption{Limiter: position error}\label{fig:tel-govep}\end{subfigure}\hfill
\begin{subfigure}[t]{0.485\linewidth}\centering\includegraphics{results/tel_gov_theta.pdf}\caption{Limiter: orientation error}\label{fig:tel-govtheta}\end{subfigure}
\caption{QP-CLF-CBF and reference-limiter data of the teleoperated trial.}
\label{fig:telqpdata}
\end{figure}

Figure~\ref{fig:tel-lcbf} is the clearest single trace of the trial's three episodes: an early right-arm engagement peaking at $220$ while it first reaches for the near cup ($t\approx9$~s), a left-arm engagement peaking at $172$ once it swings out toward the far one ($t\approx140$~s), and then the $t\approx211$--$217$~s episode, an order of magnitude higher and qualitatively different in shape --- not one spike but a repeatedly re-triggered burst, its single highest sample ($1085$) inside the broken axis's top strip, well above everything else in the trial. The shadow price is exactly zero from the trigger until $t=209.4$~s --- the row does not exist to report anything --- and every recorded solve in the $205$--$218$~s window is feasible, so each command satisfies its own tick's row; yet the margin decreased faster than the discrete counterpart of the nominal barrier condition~\eqref{eq:cbf} on most ticks of the burst, by up to $1.2$~cm per tick, the gap listed in Section~\ref{sec:concllimitations} (candidates include the velocity-dependent margin's unenforced derivative, softmin re-aggregation across pairs, penetration-depth witness normals under deep overlap, up to two ticks of barrier staleness, and the operator's own reference now fighting the row). The joint-limit shadow price of Figure~\ref{fig:tel-ljl} is comparatively quiet --- and, on its own much smaller scale, drawn with the same broken axis as Figure~\ref{fig:tel-lcbf}: a small right-arm burst ($8.5$) as it first leaves the home posture, and one sharp left-arm spike ($27$) at $t\approx140$~s as it swings out toward the far cup, with nothing in the failed-retreat window --- that episode is a collision-barrier event, not a joint-limit one. Figure~\ref{fig:tel-margin} shows the same episode from the constraint side: the left arm's safety margin sits flat near zero for the entire stall because the target pair is off the barrier entirely, then jumps in a single tick from $+2$~mm to $-3.2$~cm at $t=209.4$~s, when the bypass lifts with the gripper already inside the cup's collision envelope, before decreasing further to $-10.0$~cm at $t=214.8$~s. This reading is taken against a constraint re-armed onto an already-violated state, the direct consequence of the pair's deliberate deactivation for the grasp, and is therefore not comparable with the open-loop trial's $-2.6$~mm of Section~\ref{sec:hwswap}, recorded under a continuously enforced row. The right arm's own margin, by contrast, never drops below $-9.5$~mm, reached at $t=61.4$~s as it disengages from the failed first attempt --- a much smaller excursion than the left arm's crisis. Figure~\ref{fig:tel-dmin} carries a second, unrelated negative baseline that is not a safety event at all: raw minimum distance drops to a flat $-4.7$~cm the instant the right cup is attached ($t\approx114$~s) and stays there until the left-arm episode, because the raw metric --- unlike the barrier's own SoftMin margin above, which excludes the carried object --- scores the held cylinder against the gripper envelope it is now resting inside by construction, the same convention documented for the simulated study's clearance metrics. The further drop on top of that baseline, to $-7.5$~cm at $t=216.6$~s, coincident with the left arm's crisis, is consistent with the same left-gripper-and-blue-cup pair now unmasked --- the raw minimum begins moving exactly at $t=209.4$~s --- that is, a model overlap between the left gripper and the blue cup rather than a new contact elsewhere; the pair identity is inferred, since the bag carries no pair-identity telemetry. The velocity-dependent margin of Figure~\ref{fig:tel-dsafe} stays within its configured band throughout, $1.5$--$5.8$~cm, the same $1.5$~cm static floor as Section~\ref{sec:hwswap}; the adaptive slack weight of Figure~\ref{fig:tel-wslack} tracks each arm's own shadow price as designed, collapsing toward its floor whenever that arm's barrier is engaged and recovering toward its ceiling once it is not.

Of the reference limiter's mechanisms, the two shown here are the ones this trial exercises least: the raw and governed traces of both position error (Figure~\ref{fig:tel-govep}) and orientation error (Figure~\ref{fig:tel-govtheta}) coincide throughout for both arms, neither ever approaching the same $0.30$~m and $1.0$~rad ceilings of Section~\ref{sec:hwswap} --- the right arm's own errors, visible from $t=0$ since it alone is referenced from the start, still peak at only $10.2$~cm and $0.74$~rad, both during ordinary teleoperated motion ($t\approx71$ and $127$~s) rather than a clamp. Neither arm's tracking error ever needs governing here because a human operator's own hand motion is already smoother and slower than the quintic reference that trial drives open-loop. The limiter is not idle in this trial, however: it clamps the left arm's commanded linear and angular speed on $6.6\%$ and $11.7\%$ of ticks respectively, concentrated in the same $t\approx211$--$217$~s window --- the operator's own corrective motion during the crisis being what the limiter is actually bounding here, on a channel this figure does not show. Aggregated over the whole trial, the grasp state machine holds authority for $26.7\%$ of ticks, closely matching the $24.1\%$ obtained by summing the three grasp windows above directly from the trigger telemetry.

\begin{figure}[t]
\hwxyzlegend
\begin{subfigure}[t]{0.485\linewidth}\centering\includegraphics{results/tel_hap_pos.pdf}\caption{Handle position}\label{fig:tel-happos}\end{subfigure}\hfill
\begin{subfigure}[t]{0.485\linewidth}\centering\includegraphics{results/tel_hap_rpy.pdf}\caption{Handle orientation (roll, pitch, yaw)}\label{fig:tel-haprpy}\end{subfigure}
\vspace{0.4em}
\begin{subfigure}[t]{0.485\linewidth}\centering\includegraphics{results/tel_hap_vlin.pdf}\caption{Handle linear speed}\label{fig:tel-hapvlin}\end{subfigure}\hfill
\begin{subfigure}[t]{0.485\linewidth}\centering\includegraphics{results/tel_hap_vang.pdf}\caption{Handle angular speed}\label{fig:tel-hapvang}\end{subfigure}
\vspace{0.4em}
\begin{subfigure}[t]{0.485\linewidth}\centering\includegraphics{results/tel_hap_force.pdf}\caption{Force rendered on the handle}\label{fig:tel-hapforce}\end{subfigure}\hfill
\begin{subfigure}[t]{0.485\linewidth}\centering\includegraphics{results/tel_hap_torque.pdf}\caption{Torque rendered on the handle}\label{fig:tel-haptorque}\end{subfigure}
\caption{Haptic-device telemetry of the teleoperated trial, in the device base frame.}
\label{fig:telhaption}
\end{figure}

\textbf{Operator side.} Figure~\ref{fig:telhaption} is the same trial seen from the handle. The shaded windows are those in which the operator is in control --- deadman held and clutch released --- $87\%$ of the shown span; the white gaps are clutch holds, in which the handle is re-indexed without moving the reference (Section~\ref{sec:clutchmode}). The gaps are where the handle moves fastest: every angular-speed peak above $2$~rad/s, up to $4.1$~rad/s, falls inside one (Figure~\ref{fig:tel-hapvang}), against $2.0$~rad/s at most while in control, and the same holds for linear speed ($0.56$ against $0.36$~m/s) --- the operator's re-orientation of their own wrist between indexing strokes, never a command to the robot. The rendered force (Figure~\ref{fig:tel-hapforce}) stays within a few newtons for the whole trial, touching the $10$~N clip of the force manager exactly once, at the arm switch ($t=132.9$~s), when the tether re-anchors from the right arm's pose to the left's in a single tick. Torque (Figure~\ref{fig:tel-haptorque}) stays inside its $1$~N\,m clip (peak $0.92$); the three dense bands are not operator load at all but the vibrotactile cue rendered whenever the grasp state machine holds authority, and they coincide exactly with the three grasp windows above --- during those, no other force is impressed on the handle.

\begin{figure}[t]
\centering
\begin{subfigure}[t]{0.485\linewidth}\centering\includegraphics{results/tel_sa_alpha.pdf}\caption{Blending authority}\label{fig:tel-saalpha}\end{subfigure}\hfill
\begin{subfigure}[t]{0.485\linewidth}\centering\includegraphics{results/tel_sa_share.pdf}\caption{Policy share of the blended action}\label{fig:tel-sashare}\end{subfigure}
\vspace{0.4em}
\begin{subfigure}[t]{0.485\linewidth}\centering\includegraphics{results/tel_sa_belief.pdf}\caption{Goal belief}\label{fig:tel-sabelief}\end{subfigure}\hfill
\begin{subfigure}[t]{0.485\linewidth}\centering\includegraphics{results/tel_sa_state.pdf}\caption{Active arm and autonomous grasp}\label{fig:tel-sastate}\end{subfigure}
\caption{Shared-autonomy telemetry of the teleoperated trial.}
\label{fig:telsa}
\end{figure}

\textbf{Shared autonomy.} Figure~\ref{fig:telsa} is the arbitration layer's own account. The blending authority $\alpha$ of~\eqref{eq:blend} (Figure~\ref{fig:tel-saalpha}) drops to exactly zero throughout each grasp window, where the state machine rather than the blend owns the reference; excluding that $26.5\%$ of the trial, it is non-zero on $63\%$ of the remaining ticks, never crosses $0.73$, and averages $0.17$ --- the policy is always the minority partner. The share panel (Figure~\ref{fig:tel-sashare}) restates this in the action itself, $\alpha\|v_{\text{policy}}\|$ as a fraction of the blended twist's two contributions --- the user's share is its complement, and is not drawn --- with a median of $36\%$ where the blend is active. The belief (Figure~\ref{fig:tel-sabelief}) tracks the operator's intent through the trial's four decisions: red-side dominates the first approach, red-top takes over cleanly once the operator re-approaches from above ($t\approx100$~s), the platform goal briefly wins ($t\approx117$--$133$~s) as the operator, cup in hand, moves toward the placement station before switching arms instead, and blue-side rises to $0.998$ once the left arm is driven. Its only sustained ambiguity is blue-top against blue-side while the left arm is still far from the cup, resolved as the approach commits. Figure~\ref{fig:tel-sastate} is the discrete state under all of it: the active arm and the three windows in which a grasp, not the operator, is driving.

\section{Setup of the user study}
\label{sec:study}

The system is the apparatus of a user study run in a matched simulation environment; Section~\ref{sec:hwvalidation} reports the separate validation of the same controller and its perception pipeline on the physical robot. The design is a $2\times3$ matrix crossing the teleoperation mapping with the assistance strategy:
\begin{itemize}[nosep]
\item \textbf{Control mode}, two levels --- \textsc{clutch}, position control with indexing (Section~\ref{sec:clutchmode}), and \textsc{joystick}, spring-centred velocity control (Section~\ref{sec:joystick});
\item \textbf{Assistance strategy}, three levels --- \emph{F}, the guidance force~\eqref{eq:fguide} rendered on the handle; \emph{B}, reference-level arbitration~\eqref{eq:blend} of the user twist with the belief-weighted policy; and \emph{FB}, both together.
\end{itemize}
This yields six conditions (Table~\ref{tab:cells}). The two channels are orthogonal by construction: F acts only on the handle and never modifies the reference; B acts only on the reference and is never rendered on the handle. Consequently the combined level is the exact superposition of the two single-channel levels --- the guidance wrench under FB is identical to the one under F, and adding B changes only \emph{which} reference the base tether tracks, not the rendered assistance law. Each control mode is also implemented with neither channel active, so the system realises eight configurations in all; the two unassisted ones are used for familiarisation and as a reference trace, and are excluded from the study matrix. The two unassisted configurations are not part of the study matrix, so the study compares assistance strategies with one another and the two mappings with one another; it makes no claim that assistance outperforms unassisted teleoperation, which would require an internal baseline that this design does not run, and results reported elsewhere on other tasks cannot substitute for it. With three assistance levels (F, B, and FB) rather than the full $2\times2$ of the two channels, the independent effect of each channel and their interaction are not separately estimable: F versus FB isolates adding blending to guidance, and B versus FB isolates adding guidance to blending, and that is the extent of the channel design. The question this design does ask --- which channel, under which mapping, and how the two combine --- is not answered by that literature, so every one of the six conditions contributes to it.

\begin{table}[b]
\centering\footnotesize
\begin{tabularx}{\textwidth}{@{}lcccLL@{}}
\toprule
Condition & Code & F & B & Operator feels & Arm reference written by \\
\midrule
\textsc{clutch} guided feedback  & CF  & \checkmark & --     & tether $+$ guidance bias & teleop \\
\textsc{clutch} guided blending  & CB  & --     & \checkmark & tether (blended reference) & blending layer \\
\textsc{clutch} full guidance    & CFB & \checkmark & \checkmark & tether $+$ guidance bias & blending layer \\
\addlinespace[2pt]
\textsc{joystick} guided feedback & JF  & \checkmark & --     & spring $+$ guidance bias & teleop \\
\textsc{joystick} guided blending & JB  & --     & \checkmark & spring & blending layer \\
\textsc{joystick} full guidance   & JFB & \checkmark & \checkmark & spring $+$ guidance bias & blending layer \\
\bottomrule
\end{tabularx}
\caption{The six study conditions ($2$ control modes $\times$ $3$ assistance strategies). The code names each condition by its control mode, C for \textsc{clutch} and J for \textsc{joystick}, followed by its active assistance channels, and is the label used on every results figure. The guidance bias is~\eqref{eq:fguide}; the blending layer is the arbitration of Section~\ref{sec:blend}.}
\label{tab:cells}
\end{table}

Every condition shares Chapter~\ref{ch:method}'s architecture --- safety controller, belief estimation, goal manifolds, grasp state machine, and every shared building block (homing, damping, vibrotactile cues, authority cap) --- byte-identical throughout; only the operator-input mapping and the rendered assistance change, so a measured difference between two conditions is attributable to those factors alone.

Two asymmetries remain by design and are declared rather than equalised: the joystick base force is a \emph{homing} spring while the clutch base force is an \emph{error-tracking} tether, since each mapping mandates its own base feedback (Section~\ref{sec:columnbase}); and the guidance saturation is absolute in the clutch mapping but deadzone-exit-calibrated in the joystick mapping~\eqref{eq:kappacal}, to stay commensurate with that base-force landscape.

\subsection{Task and workspace}
\label{sec:studyworld}

Each participant works through two scenarios under every assigned condition. The two are chosen to differ in one structural respect --- whether the task \emph{requires} the arms to cooperate --- while keeping the object set, the grasp vocabulary, and the delivery-to-a-station goal common to both. This scenario factor is crossed with, and orthogonal to, the assistance factorial of Table~\ref{tab:cells}: assistance that helps a single arm run its own errand need not be the assistance that helps two arms hand an object between them.

\textbf{Relay scenario.} The task is a bimanual relay. Two industrial batteries --- modelled as plain cylinders for collision purposes --- rest on a table in front of the robot, and both must be delivered to a single station on a conveyor unit standing to the robot's right. A storage-rack shelf spans the table just above the batteries, closing the overhead approach corridor and forcing every pick and every place at the table to be a side grasp. The conveyor sits clear of the shelf, so its station admits a top approach, but only the right arm reaches it. The right arm can therefore pick its battery and carry it straight to the conveyor; the left arm cannot, and instead stages its battery on an intermediate handover disk on the table --- offset toward the robot's right so that it falls within both arms' reach --- from which the right arm relays it onward.

The handover is forced by workspace geometry alone, not by an inserted divider: only the right arm reaches the delivery station, so a run that never hands over cannot finish. Both scenes also carry decorative fixtures --- fencing, a gantry, cabinets, pallets, conveyor frames --- that give visual context but declare no collision geometry; each scenario's declared model is limited to the table, the two batteries, its own blocking structure, and one convex envelope per conveyor.

\begin{figure}[t]
\centering
\includegraphics[width=0.85\linewidth]{rackworld.png}
\caption{The relay scenario, from behind the robot: the rack shelf forces a side grasp, relayed via the handover disk (yellow) to the right arm.}
\label{fig:rackworld}
\end{figure}

\textbf{Independent scenario.} The task is two independent, concurrent pick-and-place errands, one per arm. A thin safety shield stands between the robot and the table, spanning its full width so that no lateral window is left; each arm must therefore reach \emph{over} the shield to work. Two conveyor units flank the table, mirror-symmetric about the robot's sagittal plane, each carrying its own delivery station. Each arm picks the battery on its own side and delivers it to the station on that side: the errands are symmetric, with no handover and nothing contested between the arms. Because the shield closes the low approach and leaves the overhead corridor open, every grasp here is a top grasp --- the complement of the relay scenario, where the shelf closes the overhead corridor and every grasp is a side grasp. Between them the two scenarios exercise both grasp manifolds of Section~\ref{sec:policies} and both the coordinated and the independent regimes of the bimanual controller.

Each scenario declares its placement stations as platforms --- the handover disk and the conveyor station in the relay scenario, the two mirrored conveyor stations in the independent one. Each is demandable by either arm while that arm holds a battery, and belief estimation together with guidance selects which platform is the active placement goal, exactly as they select among grasp goals (Section~\ref{sec:policies}). Like every platform, they are reference poses for the goal manifold and the guidance law alone, and none enters the collision model.

Both scenarios exercise the whole architecture, not just the driven arm. The relay's handover is a bimanual case of Section~\ref{sec:locality}'s single-arm gap: no per-arm policy decides which arm carries the battery or when to hand off, so the operator does. The head faces the same test: tracking the gripper's top-grasp approach axis (Section~\ref{sec:headtrack}) forces it above the shield's full-height collision box, a viewpoint the relay scenario never demands. Belief estimation is stressed alike: every battery carries a top and a side goal candidate regardless of which the geometry admits (Section~\ref{sec:policies}), so the belief sum~\eqref{eq:objbelief} must resolve the intended approach from poses that start close together.

\subsection{Environment model and camera use}
\label{sec:studyenv}

Two facts about the environment fix what the study tests. First, the collision world in every study trial is the declared model of Section~\ref{sec:task}, loaded from the scenario's description file; the autonomous perception pipeline of Section~\ref{sec:convergence} is not in the loop, so no trial's safety depends on the camera's estimate of the scene. Second, the head is nonetheless active throughout, running the active-arm visual servoing of Section~\ref{sec:headtrack} to keep its camera trained on the gripper being driven. That camera image is shown alongside the operator's RViz views of the session; the two are not interchangeable. The RViz views are what actually render the reference the operator is assumed throughout this thesis to be watching~(Section~\ref{sec:blend}), since a raw camera image carries no such marker. That marker is also the operator's only cue to the blending channel: it is orange while the operator's own input still dominates the blended reference, and turns green once the belief-weighted policy visibly takes over, tracking the same authority $\alpha$ of~\eqref{eq:blend} that drives the reference itself. Guidance and blending are thus felt through different senses --- the guidance bias of Section~\ref{sec:guidance} as a force on the handle, blending as a color the operator watches --- consistent with the two channels never sharing an output (Section~\ref{sec:study}). The camera image supplements this with close visual detail on the gripper itself; it serves the operator's perception, not the controller's, but it is not where the reference is seen.

\begin{figure}[t]
\centering
\includegraphics[width=0.85\linewidth]{shieldworld.png}
\caption{The independent scenario, from behind the robot: the shield forces a top grasp; each arm serves its own mirrored conveyor, independently.}
\label{fig:shieldworld}
\end{figure}

\subsection{Apparatus and instrumentation}
\label{sec:studyapparatus}

All sessions run on a single workstation --- Ubuntu 22.04 LTS, ROS~2 Humble, Gazebo Classic 11.10.2, an Intel Core i9-14900HX CPU, and an NVIDIA RTX 4090 Laptop GPU driving the Gazebo and RViz rendering. A single bringup starts the simulator and the robot model, the robot's default arm and torso controllers, the mobile-base teleoperation, and five simultaneous views of the task (Figure~\ref{fig:studyviews}). A Gazebo view from behind the robot shows the whole scene --- what an operator positioned behind the physical robot would see --- and the robot's own head camera, servoed to the driven gripper (Section~\ref{sec:headtrack}), gives a close, zoomed view of it. Three RViz views complete the operator's picture: one static view framed on the table where the batteries sit, and two views that each track one gripper, so the reference stays within at least one view's field of view whichever arm is active. All three RViz views share the Gazebo view's rear-facing azimuth --- no side or top perspective is offered --- since the device-to-robot frame convention (Section~\ref{sec:device}) is itself defined for an operator positioned behind the robot, and a perspective inconsistent with it risks disorienting a first-time participant. Each RViz view renders a filtered robot description in which every link outside the two arm-and-gripper chains is dimmed to a translucent grey, so that the operator's sight line to the hands and the task is not blocked by the robot's own torso, base, and head. The filter acts at the description source --- it rewrites the model the visualiser loads --- because per-link transparency overrides applied inside the visualiser have no visual effect on this robot's meshes. Collision and inertial data are left untouched, and the unfiltered description stays available to every other consumer.

\begin{figure}[t]
\centering
\includegraphics[width=0.85\linewidth]{generalview.png}
\caption{The general view of the operator's picture during a session.}
\label{fig:studyviews}
\end{figure}

\textbf{Operator controls.} The operator issues every discrete command through the handle's buttons, as described in Section~\ref{sec:opcontrols}. The alignment torque rendered while the clutch is held~\eqref{eq:aligntorque} is not captured by the study's translational effort measure, which is defined on the rendered force alone. The arm left inactive by a switch, and one currently executing an autonomous grasp, both bypass the adaptive slack schedule (Section~\ref{sec:sched}): their weight is pinned at the ceiling instead, so the inactive arm's hold is rigid and the grasping arm's alignment converges inside tolerance.

\textbf{Startup homing.}
\label{sec:homing}
Every session opens with a homing phase in which the only rendered wrench is a soft parking spring of the form of~\eqref{eq:fhome}, drawing the handle to a fixed neutral pose and faded in over a fixed interval, so that every participant begins from the same handle configuration; the spring is far softer than either mapping's operating spring, since the operator need not be holding the handle while it runs. The phase ends once the handle has held within position and orientation tolerances for a settle window, or after a timeout.

\subsection{Procedure}
\label{sec:studyprocedure}

Each session opens with a briefing in which the operating principles of the platform, the two control modes, and the assistance strategies the participant encounters are explained in detail. The interaction is framed explicitly as cooperative: participants are told to read the robot's guidance, whether the haptic bias or the blended reference, as collaborative assistance correcting them toward the intended goal rather than as an antagonistic force to resist, and are encouraged to work with each assistance channel rather than against it.

The briefing is followed by an unassisted familiarisation phase in a custom, simplified scene containing only the two target batteries on a bare tabletop, with every obstacle removed and neither assistance channel active. Within this scene, participants operate the haptic device with both the clutch and the joystick mapping (Sections~\ref{sec:clutchmode} and~\ref{sec:joystick}), for as long as they judge necessary. The phase serves two purposes: it lets the operator acclimatise to the device's baseline dynamics, kinematic mapping, and passive forces before either assistance channel is introduced, and it lets the operator complete the reaching task once unassisted, establishing a personal baseline against which the assistance strategies met in the formal trials can be judged. No data are recorded during familiarisation.

No individual trial carries a completion threshold or a time limit. Because a participant's data enter the analysis only as a complete six-condition record, discarding any single trial would discard the entire session rather than one data point, so nothing is gained by rushing a trial and every participant is given as long as the task requires. A complete session, from the opening briefing to the closing questionnaire, takes approximately one hour.

The six formal conditions are then run in an order fixed by a balanced Latin square, a Williams design: within each control mode, the three assistance levels are sequenced in one of the six possible orders, so that every level is preceded by every other level an equal number of times across the design, and the two three-condition blocks are further ordered so that the leading mode is balanced across participants, a balance the learning-effect analysis of the results draws on. Because three conditions admit only six distinct orderings in total, the minimal design that balances immediate order effects coincides here with the complete permutation set; crossed with the two possible leading modes it yields twelve distinct sequences, each realised by two of the twenty-four participants.

\section{Results}
\label{sec:results}

Twenty-four participants complete the full grid of six conditions on both scenarios, giving 288 analysed trials; the count follows the design, not convenience, since the schedule realises twelve counterbalanced orderings of the six conditions, each run by two participants. Every quantity below is first reduced to one number per trial, then averaged over the two scenarios so that each participant contributes exactly one value per condition, and it is those participant-level values that enter every test and every figure.

\subsection{Statistical procedure}
\label{sec:stats}

The design is fully within-subject with two factors, the control mode --- clutch or joystick --- and the assistance channel --- force guidance alone, reference blending alone, or both. Each measure is therefore examined with a two-way repeated-measures analysis of variance over mode and assistance, and, because the six conditions are also of interest individually, with a Friedman test over the six cells followed by pairwise Wilcoxon signed-rank tests whose $p$ values are corrected with the Holm step-down procedure inside each metric family. The significance level is $0.05$ throughout.

Two assumptions govern which of those two tracks carries a claim. Normality is assessed with the Shapiro--Wilk test applied to the residuals of the two-way model, $y_{pc} - \bar y_{p\cdot} - \bar y_{\cdot c} + \bar y$ with $p$ the participant and $c$ the condition, one test per measure, and to the within-participant paired differences that feed the Holm-corrected Wilcoxon tests, rather than to each condition in isolation; sphericity is assessed with Mauchly's test on the assistance factor; where sphericity is rejected, the Greenhouse--Geisser correction is applied to the degrees of freedom and the corrected $p$ value is the one reported. Sphericity is not testable, and is assumed, for a factor with only two levels, which covers the mode factor everywhere and the assistance factor in the intent-estimation and shared-authority measures that exist only where blending is active, since those are compared over the two blended assistance levels alone. Table~\ref{tab:assumptions} collects both checks together with the outcome of the three analysis-of-variance effects.

The picture that table paints is the reason the non-parametric track is treated as primary here. Residual normality fails for seven of the fourteen measures ($L$, $\bar\Phi$, $\bar\delta$, $h_{\min}$, $T_\lambda$, $\rho$, and $\Delta_{\max}$), and four fail Mauchly on the assistance factor ($L$, $\bar\delta$, $h_{\min}$, and $\bar f$ under the clutch mapping). Minimum clearance is the extreme case on both counts, for a reason that is structural rather than statistical: the barrier itself pins the measure near a fixed floor, a structural explanation the safety results take up in more detail. Consequently, the analysis of variance is reported for the structure it reveals --- whether an effect lives in the mode, in the assistance, or in their interaction, and how much variance it explains --- while every claim that one condition differs from another rests on the Holm-corrected rank tests drawn on the figures. As a check on that division of labour, a fully non-parametric factorial analysis, the aligned rank transform~\cite{wobbrock2011art}, is run on all thirty-nine tested effects; it agrees with the corrected analysis of variance at the $0.05$ level on all but one, the mode$\times$assistance interaction on the peak command alteration $\Delta_{\max}$, non-significant under the analysis of variance ($p = 0.105$) and significant under the rank-based alternative ($p = 0.031$); no claim below rests on that interaction. Both $p$ values are reported side by side in Table~\ref{tab:assumptions}. Partial eta squared, written $\eta^2_p$, accompanies each effect as its magnitude.



\subsection{Task performance}
\label{sec:resperf}

Five measures describe how the task is carried out. The task time
\begin{equation}
T = t_f - t_0,
\label{eq:mT}
\end{equation}
with $t_0$ and $t_f \in \mathbb{R}$ the first and last instants of the recording, is the coarsest of them. The distance the two hands travelled to achieve it,
\begin{equation}
L = \sum_{a \in \{R,L\}} \int_{t_0}^{t_f} \lVert \dot{p}_a(t) \rVert \, \mathrm{d}t,
\label{eq:mL}
\end{equation}
with $p_a(t) \in \mathbb{R}^3$ the position of gripper $a$, separates a direct execution from a wandering one of the same duration.

Several of the measures that follow are taken only while the operator is actually driving, and the set of such instants is fixed once here. Let $\mathcal{T}_h \subset [t_0,t_f]$ collect the instants under operator control, the autonomous grasp phases being excluded because the gripper closes on the object by design during them; let $c(t) \in \{0,1\}$ be the clutch signal, equal to one while the handle is disengaged for re-indexing, and $v_h^{\mathrm{lin}}(t) \in \mathbb{R}^3$ the linear part of the twist the operator commands. The active-driving set is
\begin{equation}
\mathcal{T}_a = \left\{ t \in \mathcal{T}_h : c(t) = 0 \ \wedge\ \lVert v_h^{\mathrm{lin}}(t) \rVert \ge v_{\text{still}} \right\},
\label{eq:Ta}
\end{equation}
with $v_{\text{still}}$ the stillness threshold below which the arbitration itself suspends (Section~\ref{sec:blend}, Appendix~\ref{app:config}), so that a hold, a re-index and a grasp contribute to no measure that is about driving.

Movement quality is summarised by the spectral arc length~\cite{balasubramanian2015sparc}. Because the measure is defined for a single movement, it is evaluated on each maximal interval $\mathcal{S}_1, \dots, \mathcal{S}_M \subset \mathcal{T}_a$ of contiguous driving that lasts at least $\tau_{\min}$, rather than on the whole recording, whose zero-velocity holds would otherwise enter the spectrum as content of their own. For segment $m$, with $\hat{V}_m(\omega) = V_m(\omega)/V_m(0)$ the magnitude spectrum of its hand speed profile $\lVert \dot{p}_A(t) \rVert$, $t \in \mathcal{S}_m$, normalised by its zero-frequency value,
\begin{equation}
\Phi_m = -\int_{0}^{\omega_c} \left[ \left(\frac{1}{\omega_c}\right)^{2} + \left(\frac{\mathrm{d}\hat{V}_m(\omega)}{\mathrm{d}\omega}\right)^{2} \right]^{1/2} \mathrm{d}\omega,
\qquad
\bar{\Phi} = \frac{1}{M} \sum_{m=1}^{M} \Phi_m,
\label{eq:mPhi}
\end{equation}
where the cutoff $\omega_c$ is the highest frequency, not above $\omega_c^{\max} = 20\pi$~rad/s, at which $\hat{V}_m(\omega)$ still reaches the amplitude threshold $\bar{V} = 0.05$, so that the arc is measured over the band that carries the movement and not over the noise floor beyond it. Because each $\Phi_m$ is the negative arc length of a spectrum, $\bar{\Phi}$ is at most zero and grows towards zero as motion becomes smoother. The mean hand speed
\begin{equation}
\bar{v} = \frac{1}{t_f-t_0} \int_{t_0}^{t_f} \lVert \dot{p}_A(t) \rVert \, \mathrm{d}t,
\label{eq:mv}
\end{equation}
taken over the arm $A$ the operator is driving, completes the kinematic picture, and the mean relaxation of the tracking constraint,
\begin{equation}
\bar{\delta} = \frac{1}{t_f-t_0} \int_{t_0}^{t_f} \delta_A(t) \, \mathrm{d}t,
\label{eq:mdelta}
\end{equation}
with $\delta_A$ the slack variable of~\eqref{eq:cost} for that arm, reports how much tracking the safety filter has to give up in order to stay feasible. A small $\bar{\delta}$ means the commanded motion is compatible with the constraints and is followed; a large one means it is not.

Figure~\ref{fig:resmain} shows these five in its panels (a) to (e). Task time and path length tell one story twice. Assistance dominates both --- it explains roughly three quarters of the variance in each, $\eta^2_p = 0.74$ for $T$ and $0.77$ for $L$ --- while the control mode explains almost none of the task time and a moderate share of the path length. Force guidance alone is the worst configuration in both modes, and adding blending to it removes about a quarter of the task time and a third of the path length. The interaction on task time is significant but small, and its shape is worth stating plainly: the two modes are not equivalent at every level of assistance. Joystick with force guidance alone is the slowest condition of the six at $306$~s, and joystick with both channels is the quickest at $196$~s, so the same mode occupies both extremes depending on what assistance accompanies it. No significant main effect of mode on task time is found, and the practice-and-order analysis gives a further reason to distrust any mode claim built on time.

Speed separates the two modes where time does not, and smoothness, once it is measured movement by movement, separates nothing. The joystick hand is the slower one, a large effect of mode with $\eta^2_p = 0.52$; the interaction on $\bar{v}$ is itself significant and large, $\eta^2_p = 0.38$, and reverses direction between modes, since assistance slows the clutch hand and speeds the joystick one. Read together with the path lengths, this says that clutch operators move faster along longer paths, which is consistent with a position-mapped device whose workspace must be re-indexed: each re-index interrupts the motion and adds distance. The segmented smoothness $\bar\Phi$ is a null result, and an instructive one. The six conditions sit between $-2.24$ and $-2.19$, no factor separates them ($p \ge 0.116$ for all three effects; Friedman $p = 0.348$, $W_{\mathrm{K}} = 0.05$), and no pair survives the Holm correction. Evaluated over the whole recording instead, holds included, the same speed profiles return a large mode effect, $\eta^2_p = 0.52$, with the joystick conditions apparently smoother by two to three units; that effect is not a property of the movements but of the pauses between them. A clutch trial contains about three times as many driving segments as a joystick one, $66$ against $34$ under force guidance alone and $51$ against $16$ with both channels, and each re-index inserts a zero-velocity hold that the spectrum of the concatenated profile reads as roughness. The individual movements are as smooth in one mapping as in the other; what the mapping changes is how many of them a trial is cut into, and that quantity is reported directly as the re-indexing count $N_c$ among the operator-effort measures.

The tracking slack is the measure on which blending acts most directly. Both factors are large and independent --- $\eta^2_p = 0.62$ for mode and $0.64$ for assistance, with no interaction --- and in both modes the two blended conditions sit roughly a quarter below the force-guidance-only one. Because $\bar\delta$ counts exactly the tracking the safety filter has to abandon, this is a statement about the commands themselves: a blended reference asks the controller for motion it is able to deliver, whereas an unassisted reference more often asks for motion that the constraints then have to refuse.



\subsection{Safety}
\label{sec:ressafety}

Safety is reported through the barrier itself rather than through a raw distance. For the arm being driven, let $h_A(t) \in \mathbb{R}$ denote the aggregated barrier value of~\eqref{eq:softmin}, that is, the smooth lower bound on the distance between that arm and the nearest obstacle, with the carried object and the grasp target excluded from the pair set. The worst clearance of a trial is
\begin{equation}
h_{\min} = \min_{t \in \mathcal{T}_h} h_A(t),
\label{eq:mhmin}
\end{equation}
with $\mathcal{T}_h$ the operator-controlled set of~\eqref{eq:Ta}. How long the filter is actually intervening is
\begin{equation}
T_{\lambda} = \int_{\mathcal{T}_h} \mathbb{1}\!\left[\lambda_A(t) > \epsilon\right] \mathrm{d}t,
\label{eq:mTlam}
\end{equation}
with $\lambda_A$ the multiplier of the collision constraint for that arm, $\mathbb{1}[\cdot]$ the indicator function, and $\epsilon = 10^{-3}$ a dual feasibility tolerance: the solver returns an exactly zero multiplier for a constraint that is not active, so any positive tolerance below the smallest genuine multiplier counts precisely the instants at which the constraint is binding, independently of how the constraint happens to be scaled. The integral runs over $\mathcal{T}_h$ for the same reason as $h_{\min}$, so that $T_{\lambda}$ measures the seconds of operator-driven motion that has to be altered to remain safe.

Panels (f) and (g) of Figure~\ref{fig:resmain} report both. The clearance result is a null result: no factor produces a significant difference in $h_{\min}$ --- mode, assistance, and their interaction all return $p > 0.26$, the Friedman test over the six conditions gives $p = 0.122$, and every condition sits between $1.9$ and $2.0$~cm. That figure is not an operator achievement but the barrier's own design margin, and the absence of any detectable difference is consistent with the safety layer, not the operator or the assistance, setting that floor, though a non-significant test does not establish equality. The severe departure from normality that Table~\ref{tab:assumptions} records for this measure has the same origin: the distribution is a narrow spike at the constraint boundary with a short tail of trials in which a genuine barrier fight pushes slightly past it, which is not a shape any test of means is designed for.

Given that the floor never moves, the informative safety measure is the second one, the time spent holding it. Here both factors are large --- $\eta^2_p = 0.43$ for mode and $0.52$ for assistance --- and the ordering is unambiguous: the filter is active for $83$~s under clutch with force guidance alone and for $52$~s under joystick with both channels, a reduction of about a third. Since no difference in attained clearance is detected, this difference is best read not as a difference in safety but in how often the operator drives into a situation from which the filter has to rescue the motion. Assistance, and blending in particular, reduces the number of such situations rather than the margin kept when one arises.



\subsection{Intent estimation and shared authority}
\label{sec:resauthority}

The remaining robot-side measures describe the assistance itself: how confidently the intention estimator commits to a goal, how well its policy agrees with the operator, how much of the motion it is granted, and how far the command it produces departs from the one it is given. The first is the time-averaged confidence in the leading goal,
\begin{equation}
\bar{b} = \frac{1}{t_f-t_0} \int_{t_0}^{t_f} \max_k b_k(t) \, \mathrm{d}t,
\label{eq:mb}
\end{equation}
with $b_k(t)$ the normalised belief score of goal $k$. The peak of that same quantity is not informative, since it reaches its ceiling in nearly every trial of every condition; the time average is what distinguishes an estimator that commits early and holds from one that keeps changing its mind. Agreement between operator and policy is measured as the arbitration itself measures it (Section~\ref{sec:blend}), one channel at a time, since a cosine between quantities in metres per second and one between quantities in radians per second cannot be combined in a single inner product. With $(v_h^{\mathrm{lin}}, \omega_h)$ the linear and angular parts of the twist the operator commands and $(\pi^{\mathrm{lin}}, \omega_\pi)$ those of the belief-weighted policy twist $\bar{\pi}$ of~\eqref{eq:blend},
\begin{equation}
\rho_{\mathrm{lin}}(t) = \frac{v_h^{\mathrm{lin}}(t)^{\top} \pi^{\mathrm{lin}}(t)}{\lVert v_h^{\mathrm{lin}}(t)\rVert \, \lVert \pi^{\mathrm{lin}}(t)\rVert}, \qquad
\rho_{\mathrm{ang}}(t) = \frac{\omega_h(t)^{\top} \omega_\pi(t)}{\lVert \omega_h(t)\rVert \, \lVert \omega_\pi(t)\rVert},
\label{eq:mrhoc}
\end{equation}
and the instantaneous agreement is the mean over the channels the operator is driving, $\mathcal{C}_{\mathrm{act}}(t) \subseteq \{\mathrm{lin}, \mathrm{ang}\}$, the linear one on all of $\mathcal{T}_a$ and the angular one where $\lVert \omega_h(t) \rVert \ge \omega_{\text{still}}$,
\begin{equation}
\rho(t) = \frac{1}{|\mathcal{C}_{\mathrm{act}}(t)|} \sum_{c \in \mathcal{C}_{\mathrm{act}}(t)} \rho_c(t), \qquad
\rho = \frac{1}{|\mathcal{T}_a|} \sum_{t \in \mathcal{T}_a} \rho(t),
\label{eq:mrho}
\end{equation}
with a channel's cosine taken as zero where the policy is silent in it. The share of the driving during which the policy held the larger part of the authority is
\begin{equation}
\phi_\alpha = \frac{\left|\left\{ t \in \mathcal{T}_a : \alpha(t) > 1/2 \right\}\right|}{\left|\mathcal{T}_a\right|},
\label{eq:mphi}
\end{equation}
with $\alpha(t) \in [0,1]$ the arbitration weight of~\eqref{eq:blend}. The ratio is taken over $\mathcal{T}_a$ and not over the trial because the arbitration suspends itself, with $\alpha = 0$, while the clutch is held, while the hand is still, and throughout the autonomous grasp phases: a share of the whole recording would measure how long the operator paused, and would do so differently in the two mappings, rather than who led while they drove. Finally, the extent to which arbitration rewrote the operator's command is
\begin{equation}
\bar{\Delta} = \frac{1}{|\mathcal{T}_a|} \sum_{t \in \mathcal{T}_a} \lVert v_b(t) - v_h(t) \rVert, \qquad
\Delta_{\max} = \max_{t \in \mathcal{T}_a} \lVert v_b(t) - v_h(t) \rVert,
\label{eq:mDelta}
\end{equation}
with $v_b(t)$ the blended twist finally sent to the safety filter and the norms taken over the linear part alone, so that both are expressed in metres per second. These two are the transparency measure of the shared-control literature: a system that executes faithfully what it is asked keeps them near zero, and one that overrides its operator does not. Neither is given a preferred direction here, because altering the command is the declared purpose of the blending channel; whether a given amount of alteration is help or interference is decided by the measures that surround it, not by $\bar{\Delta}$ alone. The four measures of this section exist only where arbitration exists, that is, in the four conditions that include the blending channel.

Panel (h) of Figure~\ref{fig:resmain} and panels (a) to (d) of Figure~\ref{fig:resextra} report them. Confidence is a mode effect and nothing else: the estimator is surer under joystick control, $\eta^2_p = 0.60$, and entirely indifferent to which assistance accompanies it. The mechanism is the one built into the estimator, which scores the operator's twist against goal-directed policies; a velocity mapping produces a continuous stream of such evidence, while a position mapping interrupts it at every re-index. Agreement and authority then move together and, unusually for this study, both show a significant interaction, $\eta^2_p = 0.30$ and $0.45$. Agreement roughly doubles from the weakest condition to the strongest, from $0.22$ under clutch with blending to $0.47$ under joystick with both channels, and the share of driving time in which the policy leads grows from $0.22$ to $0.53$; in both, the gain from adding force guidance to blending is several times larger under joystick control than under clutch, $0.11$ against $0.03$ in agreement and $0.13$ against $0.01$ in led time. The two channels of the agreement tell the same story separately: the linear channel agrees more than the angular one in every condition, and it is the angular channel that gains most from the guidance force under the joystick, from $0.27$ to $0.43$, against $0.13$ to $0.16$ under the clutch. This is the clearest mechanistic finding of the study: force guidance and reference blending act through the same channel in a velocity mapping, since the force the operator feels is a velocity cue and the blending alters a velocity, whereas in a position mapping the two speak different languages and reinforce each other less.

The two transparency measures are the ones that most repay careful reading, because they disagree. On the average, arbitration alters the joystick command more --- $\bar{\Delta}$ rises from $1.7$~cm\,s$^{-1}$ under clutch with blending to $2.4$~cm\,s$^{-1}$ under joystick with both channels. On the worst single instant, the ordering reverses and the margin widens: $\Delta_{\max}$ reaches $19$~cm\,s$^{-1}$ under clutch with both channels against $8$~cm\,s$^{-1}$ under joystick. Arbitration in the joystick mapping is therefore a steady, small hand on the command, while in the clutch mapping it mostly stands aside and then intervenes in large discrete corrections. The same behaviour appears in the peak measure's departure from normality in Table~\ref{tab:assumptions}, which is produced by exactly those rare large interventions. Taken with the agreement result, this explains why the joystick conditions score better without the operator being overridden more violently: the assistance is present more of the time and disagrees with the operator less when it is.



\subsection{Operator effort at the interface}
\label{sec:reseffort}

The measures considered so far are recorded on the robot side. What the operator physically does is described by the wrench the device renders against the hand,
\begin{equation}
\bar{f} = \frac{1}{t_f-t_0} \int_{t_0}^{t_f} \lVert f_d(t) \rVert \, \mathrm{d}t,
\label{eq:mf}
\end{equation}
with $f_d(t) \in \mathbb{R}^3$ the commanded device force, and, in the clutch mapping only, by the number of re-indexing operations
\begin{equation}
N_c = \left| \left\{ t : c(t^-) = 0 \ \wedge \ c(t^+) = 1 \right\} \right|,
\label{eq:mNc}
\end{equation}
that is, the count of rising edges of the clutch signal $c(t) \in \{0,1\}$.

These two must be read within a mode and never across one. The two mappings render different things through the same actuator: in the clutch mapping the wrench is a tether that pulls the handle towards the gripper pose, whereas in the joystick mapping it is a centring spring that returns the handle to neutral and whose magnitude therefore encodes commanded velocity rather than tracking error. Both are newtons, but a clutch newton and a joystick newton do not denote the same physical situation, and the absolute levels in panels (e) and (f) of Figure~\ref{fig:resextra} are consequently not comparable between them. Re-indexing has no counterpart at all outside the clutch mapping. For this reason no mode effect is reported for either measure in Table~\ref{tab:assumptions}, and each is analysed as a one-way repeated-measures comparison of the three assistance levels inside its own mode.

Within the clutch mapping the effect is large and unambiguous. Mean rendered force falls from $0.82$~N under force guidance alone to $0.48$~N with blending, a reduction of about $40$~per cent with $\eta^2_p = 0.77$, and the number of re-indexing operations falls from $22$ to $16$ per trial with $\eta^2_p = 0.44$. Blending therefore reduces both the force the operator absorbs and the number of times the workspace limit must be worked around. Within the joystick mapping the same comparison returns nothing: the three levels sit within $0.1$~N of each other and the test does not separate them, $p = 0.193$. This asymmetry follows from what each wrench represents. The clutch tether grows with the discrepancy between the handle and the gripper, and blending reduces that discrepancy by steering the reference towards the goal, so the tether relaxes; the joystick centring spring is a function of how far the operator displaces the handle, which is a matter of how fast they choose to drive and is largely unaffected by what the arbitration does downstream.







\begin{figure}[p]
\centering
\begin{subfigure}[t]{0.485\linewidth}\centering\includegraphics{results/res_T.pdf}\caption{Task time}\label{fig:resT}\end{subfigure}\hfill
\begin{subfigure}[t]{0.485\linewidth}\centering\includegraphics{results/res_L.pdf}\caption{Hand path length}\label{fig:resL}\end{subfigure}
\vspace{0.35em}
\begin{subfigure}[t]{0.485\linewidth}\centering\includegraphics[width=\linewidth]{results/res_Phi.pdf}\caption{Movement smoothness}\label{fig:resPhi}\end{subfigure}\hfill
\begin{subfigure}[t]{0.485\linewidth}\centering\includegraphics[width=\linewidth]{results/res_vbar.pdf}\caption{Mean hand speed}\label{fig:resv}\end{subfigure}
\vspace{0.35em}
\begin{subfigure}[t]{0.485\linewidth}\centering\includegraphics[width=\linewidth]{results/res_deltabar.pdf}\caption{Mean tracking slack}\label{fig:resdelta}\end{subfigure}\hfill
\begin{subfigure}[t]{0.485\linewidth}\centering\includegraphics[width=\linewidth]{results/res_hmin.pdf}\caption{Minimum clearance}\label{fig:reshmin}\end{subfigure}
\vspace{0.35em}
\begin{subfigure}[t]{0.485\linewidth}\centering\includegraphics{results/res_Tlambda.pdf}\caption{Filter active time}\label{fig:resTlam}\end{subfigure}\hfill
\begin{subfigure}[t]{0.485\linewidth}\centering\includegraphics[width=\linewidth]{results/res_bbar.pdf}\caption{Mean intent confidence}\label{fig:resb}\end{subfigure}
\caption{Condition means with the standard error of the mean; brackets mark the Holm-corrected significant pairs. Symbols and units are defined in Sections~\ref{sec:resperf}--\ref{sec:resauthority}.}
\label{fig:resmain}
\end{figure}

\begin{figure}[p]
\centering
\begin{subfigure}[t]{0.485\linewidth}\centering\includegraphics{results/res_rho.pdf}\caption{Operator--policy agreement}\label{fig:resrho}\end{subfigure}\hfill
\begin{subfigure}[t]{0.485\linewidth}\centering\includegraphics{results/res_phialpha.pdf}\caption{Autonomy-led time}\label{fig:resphi}\end{subfigure}
\vspace{0.35em}
\begin{subfigure}[t]{0.485\linewidth}\centering\includegraphics[width=\linewidth]{results/res_Dmean.pdf}\caption{Mean command alteration}\label{fig:resDmean}\end{subfigure}\hfill
\begin{subfigure}[t]{0.485\linewidth}\centering\includegraphics[width=\linewidth]{results/res_Dpeak.pdf}\caption{Peak command alteration}\label{fig:resDpeak}\end{subfigure}
\vspace{0.35em}
\begin{subfigure}[t]{0.485\linewidth}\centering\includegraphics{results/res_fbarC.pdf}\caption{Rendered force, clutch}\label{fig:resfC}\end{subfigure}\hfill
\begin{subfigure}[t]{0.485\linewidth}\centering\includegraphics[width=\linewidth]{results/res_fbarJ.pdf}\caption{Rendered force, joystick}\label{fig:resfJ}\end{subfigure}
\vspace{0.35em}
\begin{subfigure}[t]{0.485\linewidth}\centering\includegraphics{results/res_Nc.pdf}\caption{Re-indexing operations}\label{fig:resNc}\end{subfigure}\hfill
\begin{subfigure}[t]{0.485\linewidth}\centering\includegraphics{results/res_learn_slot.pdf}\caption{Score along the session}\label{fig:reslearnslot}\end{subfigure}
\caption{Panels (a) to (d) exist only where blending is active and show four conditions; (e) to (g) compare assistance inside one mapping and their levels are not comparable across panels; (h) tracks three standardised scores over the session.}
\label{fig:resextra}
\end{figure}

\subsection{Learning effect}
\label{sec:reslearning}

The counterbalanced order of Section~\ref{sec:studyprocedure} balances the mode averages but not practice, so two checks test whether learning contaminates the comparisons above, the first reported in panel (h) of Figure~\ref{fig:resextra}.

The first asks whether performance drifts along the session. The composite score, obtained by averaging the standardised family scores, is flat: its slope over the six slots is not significant, so the study as a whole shows no overall practice trend. Two families do drift, in opposite directions. Time and effectiveness improves, which is the expected shape of learning, while operator effort deteriorates, that is, participants work the handle harder by the end of the session than at its beginning. Fatigue is the natural reading of a measure that worsens monotonically while task performance improves, and it is a caution for any effort claim that relies on ordering.

The second check asks whether the difference between the two modes depends on which of them a participant meets first. Splitting the twenty-four participants into the twelve who meet the clutch mapping first and the twelve who meet the joystick mapping first, the joystick-minus-clutch difference in composite score is statistically indistinguishable between the two groups, $p = 0.544$ by rank-sum test, the two group means differing by less than a tenth of a standard deviation. This is the single most important number for the mode conclusion, because it establishes that the advantage is not an artefact of practice accumulated before the second mode is reached. The same check fails, however, for the time-and-effectiveness and operator-effort families, which is why no claim in Section~\ref{sec:resperf} rests on task time alone, and why the speed comparison between modes is reported there as undetermined rather than as a result.






\begin{table}[p]
\centering
\footnotesize
\caption[Assumption checks and analysis-of-variance outcomes for every reported measure]{Assumption checks and analysis-of-variance outcomes for every reported measure, one line per tested effect. Symbols are defined with the corresponding measure in Sections~\ref{sec:resperf}--\ref{sec:reseffort}; the value in parentheses after each symbol is the $p$ value of the Shapiro--Wilk test on the residuals of the two-way model, one test per measure. Mauchly's $p$ and the Greenhouse--Geisser $\varepsilon$ are those of the effect on the line, the interaction carrying its own; where Mauchly rejects sphericity at $0.05$ the degrees of freedom of $F$ are the adjusted ones, $\varepsilon\,\mathrm{df}$, and the $p$ value is the corrected one. A dash marks a factor with two levels, for which sphericity holds by construction. The last column is the $p$ value of the aligned-rank-transform analysis of the same effect~\cite{wobbrock2011art}, or the Friedman $p$ value for the one-way measures $\bar f$ and $N_c$, which are compared within one mapping only.}
\label{tab:assumptions}
\setlength{\tabcolsep}{3.5pt}\scriptsize
\begin{tabular}{@{}llcccccc@{}}
\toprule
Measure (SW $p$) & Effect & Mauchly $p$ & $\varepsilon$ & $F(\mathrm{df}_1,\mathrm{df}_2)$ & $p$ & $\eta^2_p$ & ART / Friedman $p$ \\
\midrule
$T$ ($0.415$) & Mode & --- & --- & $F(1,23)=0.93$ & $0.345$ & $0.04$ & $0.169$ \\
 & Assist. & $0.180$ & $0.87$ & $F(2,46)=64.53$ & $<0.001$ & $0.74$ & $<0.001$ \\
 & Mode$\times$Assist. & $0.007$ & $0.73$ & $F(1.47,33.80)=4.23$ & $0.033$ & $0.16$ & $0.007$ \\
\addlinespace[2pt]
$L$ ($<0.001$) & Mode & --- & --- & $F(1,23)=12.31$ & $0.002$ & $0.35$ & $<0.001$ \\
 & Assist. & $<0.001$ & $0.67$ & $F(1.34,30.73)=75.84$ & $<0.001$ & $0.77$ & $<0.001$ \\
 & Mode$\times$Assist. & $0.008$ & $0.74$ & $F(1.48,34.00)=2.78$ & $0.090$ & $0.11$ & $0.099$ \\
\addlinespace[2pt]
$\bar{\Phi}$ ($0.004$) & Mode & --- & --- & $F(1,23)=2.66$ & $0.116$ & $0.10$ & $0.276$ \\
 & Assist. & $0.165$ & $0.87$ & $F(2,46)=0.03$ & $0.969$ & $0.00$ & $0.870$ \\
 & Mode$\times$Assist. & $0.029$ & $0.78$ & $F(1.57,36.06)=1.40$ & $0.257$ & $0.06$ & $0.294$ \\
\addlinespace[2pt]
$\bar{v}$ ($0.217$) & Mode & --- & --- & $F(1,23)=25.24$ & $<0.001$ & $0.52$ & $<0.001$ \\
 & Assist. & $0.389$ & $0.92$ & $F(2,46)=2.83$ & $0.070$ & $0.11$ & $0.129$ \\
 & Mode$\times$Assist. & $0.182$ & $0.87$ & $F(2,46)=14.23$ & $<0.001$ & $0.38$ & $<0.001$ \\
\addlinespace[2pt]
$\bar{\delta}$ ($<0.001$) & Mode & --- & --- & $F(1,23)=37.21$ & $<0.001$ & $0.62$ & $<0.001$ \\
 & Assist. & $0.005$ & $0.73$ & $F(1.45,33.40)=40.01$ & $<0.001$ & $0.64$ & $<0.001$ \\
 & Mode$\times$Assist. & $0.010$ & $0.75$ & $F(1.49,34.34)=0.26$ & $0.707$ & $0.01$ & $0.555$ \\
\addlinespace[2pt]
$h_{\min}$ ($<0.001$) & Mode & --- & --- & $F(1,23)=0.00$ & $0.983$ & $0.00$ & $0.640$ \\
 & Assist. & $<0.001$ & $0.68$ & $F(1.35,31.13)=1.36$ & $0.262$ & $0.06$ & $0.091$ \\
 & Mode$\times$Assist. & $<0.001$ & $0.60$ & $F(1.20,27.69)=0.18$ & $0.722$ & $0.01$ & $0.440$ \\
\addlinespace[2pt]
$T_{\lambda}$ ($<0.001$) & Mode & --- & --- & $F(1,23)=17.20$ & $<0.001$ & $0.43$ & $<0.001$ \\
 & Assist. & $0.069$ & $0.82$ & $F(2,46)=25.32$ & $<0.001$ & $0.52$ & $<0.001$ \\
 & Mode$\times$Assist. & $0.015$ & $0.76$ & $F(1.52,34.93)=0.82$ & $0.420$ & $0.03$ & $0.398$ \\
\addlinespace[2pt]
$\bar{b}$ ($0.131$) & Mode & --- & --- & $F(1,23)=34.83$ & $<0.001$ & $0.60$ & $<0.001$ \\
 & Assist. & $0.152$ & $0.86$ & $F(2,46)=1.73$ & $0.189$ & $0.07$ & $0.237$ \\
 & Mode$\times$Assist. & $0.550$ & $0.95$ & $F(2,46)=1.83$ & $0.172$ & $0.07$ & $0.216$ \\
\addlinespace[2pt]
$\rho$ ($0.011$) & Mode & --- & --- & $F(1,23)=167.6$ & $<0.001$ & $0.88$ & $<0.001$ \\
 & Assist. & --- & --- & $F(1,23)=20.29$ & $<0.001$ & $0.47$ & $<0.001$ \\
 & Mode$\times$Assist. & --- & --- & $F(1,23)=10.09$ & $0.004$ & $0.30$ & $<0.001$ \\
\addlinespace[2pt]
$\phi_\alpha$ ($0.851$) & Mode & --- & --- & $F(1,23)=256.2$ & $<0.001$ & $0.92$ & $<0.001$ \\
 & Assist. & --- & --- & $F(1,23)=20.21$ & $<0.001$ & $0.47$ & $<0.001$ \\
 & Mode$\times$Assist. & --- & --- & $F(1,23)=18.70$ & $<0.001$ & $0.45$ & $<0.001$ \\
\addlinespace[2pt]
$\bar{\Delta}$ ($0.432$) & Mode & --- & --- & $F(1,23)=29.93$ & $<0.001$ & $0.57$ & $<0.001$ \\
 & Assist. & --- & --- & $F(1,23)=10.40$ & $0.004$ & $0.31$ & $0.003$ \\
 & Mode$\times$Assist. & --- & --- & $F(1,23)=0.04$ & $0.835$ & $0.00$ & $0.489$ \\
\addlinespace[2pt]
$\Delta_{\max}$ ($<0.001$) & Mode & --- & --- & $F(1,23)=36.37$ & $<0.001$ & $0.61$ & $<0.001$ \\
 & Assist. & --- & --- & $F(1,23)=4.62$ & $0.042$ & $0.17$ & $0.014$ \\
 & Mode$\times$Assist. & --- & --- & $F(1,23)=2.84$ & $0.105$ & $0.11$ & $0.031$ \\
\addlinespace[2pt]
$\bar{f}$ (clutch) ($0.477$) & Assist. & $<0.001$ & $0.66$ & $F(1.31,30.14)=75.70$ & $<0.001$ & $0.77$ & $<0.001$ \\
\addlinespace[2pt]
$\bar{f}$ (joystick) ($0.439$) & Assist. & $0.160$ & $0.87$ & $F(2,46)=1.71$ & $0.193$ & $0.07$ & $0.311$ \\
\addlinespace[2pt]
$N_c$ ($0.917$) & Assist. & $0.592$ & $0.96$ & $F(2,46)=17.82$ & $<0.001$ & $0.44$ & $<0.001$ \\
\bottomrule
\end{tabular}
\end{table}

\FloatBarrier
\subsection{Subjective experience}
\label{sec:ressubj}

Twenty-four volunteers take part, eighteen men and five women, one of whom does not return the opening form. Ages range from $20$ to $58$~years, with a mean of $28.0$ and a standard deviation of $7.7$ over the twenty-two who report one. Twenty are right-handed and three report no strong hand preference. Self-rated familiarity with teleoperation devices averages $2.3$ on a five-point scale, with eleven participants placing themselves at the lowest level and six in the top two, so the sample is predominantly inexperienced. Every participant returns all six condition forms and the closing form, which yields 144 condition reports with no missing rating and forty-eight complete preference rankings.

The questionnaire is administered once per condition, after the participant has attempted that condition on both scenarios, and carries six items answered on a seven-point ordinal scale: three workload items drawn from the NASA task load index~\cite{hart1988nasatlx} and three human--robot-interaction items, of which the trust item follows Jian, Bisantz, and Drury~\cite{jian2000trust}. Mental demand $S_{\mathrm{md}}$ and physical demand $S_{\mathrm{pd}}$ are costs, on which a high score is unfavourable. Perceived success $S_{\mathrm{pf}}$, the sense of being in control of the robot's motion $S_{\mathrm{ct}}$, trust that the robot does what is intended $S_{\mathrm{tr}}$, and the judgement that the motion is smooth and comfortable $S_{\mathrm{sm}}$ are benefits, on which a high score is favourable. Reversing the two demand items as $8-S$ and averaging all six yields one subjective quality index $\bar{S}$ per condition, on the same seven-point scale and oriented so that a higher value denotes a better experience; the items are consistent enough for that summary to carry meaning, with Cronbach's $\alpha = 0.853$ over the 144 reports. Because the responses are ordinal, they are examined with the rank procedure of Section~\ref{sec:stats} alone, and Kendall's coefficient of concordance $W_{\mathrm{K}}$ accompanies each Friedman test as the magnitude of its effect.

\begin{figure}[t]
\centering
\begin{subfigure}[t]{0.32\linewidth}\centering\includegraphics[width=\linewidth]{results/res_subj_md.pdf}\caption{Mental demand}\label{fig:subjmd}\end{subfigure}\hfill
\begin{subfigure}[t]{0.32\linewidth}\centering\includegraphics[width=\linewidth]{results/res_subj_pd.pdf}\caption{Physical demand}\label{fig:subjpd}\end{subfigure}\hfill
\begin{subfigure}[t]{0.32\linewidth}\centering\includegraphics[width=\linewidth]{results/res_subj_pf.pdf}\caption{Perceived success}\label{fig:subjpf}\end{subfigure}
\vspace{0.35em}
\begin{subfigure}[t]{0.32\linewidth}\centering\includegraphics[width=\linewidth]{results/res_subj_ct.pdf}\caption{Sense of control}\label{fig:subjct}\end{subfigure}\hfill
\begin{subfigure}[t]{0.32\linewidth}\centering\includegraphics[width=\linewidth]{results/res_subj_tr.pdf}\caption{Trust}\label{fig:subjtr}\end{subfigure}\hfill
\begin{subfigure}[t]{0.32\linewidth}\centering\includegraphics[width=\linewidth]{results/res_subj_sm.pdf}\caption{Motion comfort}\label{fig:subjsm}\end{subfigure}
\caption{Questionnaire responses by condition, as means over the twenty-four participants with the standard error of the mean. Every item is answered on the same scale, running from $1$ at the axis floor to $7$; the vertical rule separates the clutch mapping from the joystick mapping, and the conditions are coded as in Table~\ref{tab:cells}. Brackets above the scale give the Holm-corrected Wilcoxon $p$ values of the two contrasts against force guidance alone within each mapping. Symbols are defined in the text.}
\label{fig:ressubj}
\end{figure}

\begin{table}[t]
\centering
\small
\caption{Effect of the assistance channel on each questionnaire item, collapsed over the two mappings, under force guidance alone (F), reference blending alone (B), and both together (FB). The Friedman test runs over those three settings, with Kendall's coefficient of concordance $W_{\mathrm{K}}$ as its effect size, and the last three columns are Holm-corrected Wilcoxon $p$ values. $\bar{S}$ is the composite index; a high score is unfavourable for $S_{\mathrm{md}}$ and $S_{\mathrm{pd}}$ and favourable otherwise. Symbols are defined in Section~\ref{sec:ressubj}.}
\label{tab:subjassist}
\begin{tabular}{lcccccccccc}
\toprule
 & \multicolumn{3}{c}{Mean rating} & \multicolumn{3}{c}{Friedman} & \multicolumn{3}{c}{Holm-corrected $p$} \\
\cmidrule(lr){2-4}\cmidrule(lr){5-7}\cmidrule(lr){8-10}
 & F & B & FB & $\chi^2(2)$ & $p$ & $W_{\mathrm{K}}$ & F--B & F--FB & B--FB \\
\midrule
$S_{\mathrm{md}}$ & $4.35$ & $3.23$ & $3.21$ & $19.50$ & $<0.001$ & $0.41$ & $0.002$   & $0.002$   & $0.985$ \\
$S_{\mathrm{pd}}$ & $3.62$ & $2.62$ & $2.90$ & $18.30$ & $<0.001$ & $0.38$ & $0.008$   & $0.009$   & $0.062$ \\
$S_{\mathrm{pf}}$ & $5.17$ & $5.90$ & $6.08$ & $17.26$ & $<0.001$ & $0.36$ & $0.003$   & $0.004$   & $0.154$ \\
$S_{\mathrm{ct}}$ & $4.54$ & $5.77$ & $5.65$ & $17.46$ & $<0.001$ & $0.36$ & $<0.001$  & $0.008$   & $0.818$ \\
$S_{\mathrm{tr}}$ & $4.42$ & $5.81$ & $5.77$ & $24.21$ & $<0.001$ & $0.50$ & $<0.001$  & $<0.001$  & $0.873$ \\
$S_{\mathrm{sm}}$ & $4.29$ & $5.65$ & $5.83$ & $22.07$ & $<0.001$ & $0.46$ & $0.001$   & $<0.001$  & $0.180$ \\
\midrule
$\bar{S}$         & $4.41$ & $5.55$ & $5.54$ & $27.85$ & $<0.001$ & $0.58$ & $<0.001$  & $<0.001$  & $0.987$ \\
\bottomrule
\end{tabular}
\end{table}

\begin{figure}[t]
\centering
\includegraphics{results/res_subj_ranks.pdf}
\caption{Preference ranking of the three assistance settings within each mapping, as the number of the twenty-four participants placing each condition first, second, or third. The horizontal axis counts participants, so every row sums to twenty-four, and every segment is annotated with its own count. The horizontal rule separates the clutch mapping from the joystick mapping, and the conditions are coded as in Table~\ref{tab:cells}.}
\label{fig:subjranks}
\end{figure}

Figure~\ref{fig:ressubj} reports every item in every condition, and one pattern governs all six. The assistance channel separates the conditions, it separates them in the same direction under both mappings, and force guidance alone is the least favourable of the three settings on every item. Table~\ref{tab:subjassist} gives the statistics collapsed over the mapping. The Friedman test rejects for all six items and for the index, with $W_{\mathrm{K}}$ between $0.36$ and $0.58$, a large effect for a within-subject design of this size, and the index rises from $4.41$ under force guidance alone to $5.55$ under blending and $5.54$ under both channels together.

Two features of that pattern matter more than the omnibus result. The first is which items move most. Blending removes authority from the operator, yet the two items it raises most are precisely the ones that a loss of authority might have been expected to cost: the sense of control climbs from $4.54$ to $5.77$, and trust from $4.42$ to $5.81$. Participants do not experience the arbitration as authority taken away but as the robot doing what they mean. The second is what does not move. Blending alone and blending with force guidance are indistinguishable on every item and on the index, the corrected $p$ never falling below $0.06$, so no evidence is found that the haptic channel adds anything the operator can detect once the reference is already being blended. The one exception is physical demand inside the clutch mapping, where the two blended settings do differ, $p = 0.026$; it has no counterpart under the joystick mapping and does not survive collapsing the two.

No subjective difference between the two mappings is detected. None of the six items separates the two, and neither does the index, which differs by $0.13$ points in favour of the clutch mapping, $p = 0.498$. The largest of the seven contrasts is motion comfort, at $0.29$ points, and it too favours the clutch mapping, $p = 0.092$. This deserves to be stated against the objective record, because the two disagree: the joystick mapping produces the shorter hand path $L$ and the more confident intent estimate $\bar{b}$, yet none of that advantage appears in what the operators report, and the one subjective trend that does appear points the other way.



Ranking the three assistance settings inside each mapping reproduces that result through a different question format. Figure~\ref{fig:subjranks} gives the distributions. Force guidance alone is placed last by seventeen of the twenty-four participants under the clutch mapping and by twenty-one of twenty-four under the joystick mapping, where it draws no first place at all. The two blended settings are again indistinguishable from one another, $p = 0.663$ in both mappings, so the ordering the ratings establish is confirmed by an instrument that forces a choice rather than inviting a score.

Asked at the end of the session which of the two mappings is the better overall, thirteen participants choose the clutch and eleven the joystick, an even split, exact binomial $p = 0.839$. The question offers no neutral option, so the count cannot separate a divided population from an indifferent one; read together with the flat index contrast above, indifference is at least as plausible a reading as division.

%───────────────────────────────────────────────────────────────────────────────
